\documentclass{article}
\usepackage[utf8]{inputenc}
\usepackage[russian]{babel} % Для поддержки русского языка

\title{Пример научной статьи}
\author{Автор}
\date{\today}

\begin{document}

\maketitle

\section{Введение}
Это вводная часть статьи. Здесь обычно описываются цели работы, актуальность темы и краткий обзор существующих решений. Введение должно заинтересовать читателя и дать общее представление о содержании статьи.

\section{Основная часть}
В основной части статьи подробно рассматриваются методы, результаты и обсуждения. Эта часть может содержать несколько подразделов, таблицы, графики и формулы. В данном примере основная часть intentionally оставлена пустой для краткости.

\section{Заключение}
В заключении подводятся итоги работы. Здесь формулируются основные выводы, описываются достигнутые результаты и могут быть намечены перспективы дальнейших исследований. Заключение должно логично завершать статью и отвечать на вопросы, поставленные во введении.

\end{document}